\section{Complete classification of particles and physical states}
\label{sec:catalogo-realizaciones-tipadas-rev8}
\index{physical sectors!typed inventory}
\index{mass atlas!81/56/13/324}

The realizations are ordered by object, multiplicity, and codomain. The mathematical domain contains \(81\) positions, grouped into
distinct codomains. The mathematical domain contains \(81\) positions, grouped into
\(56\) route families and \(13\) extension classes; the \(324\)
oriented routes adjoin the initial orientation and do not represent additional masses
or particles. The ``evaluation operator'' column does not impose a mass where the object requires
a null chart, a prior composition, a topological invariant, or a
prolongation obstruction.

\subsection{Spectral definitions of particle, spin, color, and flavor}

A \emph{particle} is a surviving section realized as an oriented route, endowed with a register, a symmetry representation, a spectral operator, and a physical chart. A fermion realizes the exterior grading and the CAR relations; a boson realizes the symmetric completion and the CCR relations. Spin records the holonomy of orientation: the return of \(2\pi\) closes the bosonic phase, whereas the spinorial section recovers its state after \(4\pi\). This distinction links the double areal--volumetric sheet to parity and Pauli exclusion.

The \emph{color} is the ternary representation that lifts the coordinate
\(r-c\pmod3\) to the fundamental orbit and thereafter to the adjoint of
\(\mathfrak{su}_3\). The \emph{flavor} identifies a realization basis
within a representation with fixed charge, color, and spin; CKM and PMNS
transport between the interaction basis and the spectral basis. Neutrinos
are neutral fermionic sections: their flavor labels belong to the
mixing reader and their masses to the eigenvectors of the corresponding operator.

A virtual particle is a finite section with a ledger and propagator up to
a horizon of prolongation; a Majorana realization is the real
fermionic sector fixed by conjugation; and the Fibonacci, Ising, and
\(\mathbb Z/6\mathbb Z\) sectors publish quantum dimensions, fusion rules, and
braid characters. Each definition thus preserves the codomain assigned to it by the
structure.

